% ============================================================
% 第十一章 · 最终得分
% 对应大纲：# 十一、最终得分
%
% 本章是写给任课教师的评分申请，读的人不是组内成员，措辞以陈述事实为主，
% 不写组内才知道的简称与玩笑。
%
% 事实来源：
%   chapters/05-management.tex   版本时间轴、229 次提交与 73 个 PR 的统计
%   chapters/02-game-design.tex  表 2.3 的项目分工
%   chapters/09-individual-summary.tex  各人自述与 PR 记录
%   第一章、第三章               剧情节点数、转场数、视频数、场景数
%
% 文中数字与上述章节保持一致，改动任一数字时需同步核对第五章末段与第二章
% 分工表。日期区间用 -- 连接，不写「9.3 至 9.8」。
%
% 组员顺序与 tex/info.tex 的成员名单、第九章保持一致。
% ============================================================

\chapter{最终得分}

《\ProjectName》从 9 月 2 日立项、9 月 20 日交稿，前后十八天。这十八天里，小组在 \texttt{main} 上\textbf{累计提交 229 次，开出 73 个合并请求，其中 68 个合入。交付的成品包含 18 个探索场景、4 个小游戏、5 个结局、11 处转场效果和 7 段过场视频，中间经历过难以计次大量合作，对接，评测与联调}。最终游戏从入口页进入可以一路玩到终局，中间不需要任何跳过，也不需要打补丁，基本达到了最初的项目结项要求。

四个版本是首尾相接的：上一代还在收尾，下一代的需求评审已经开完。这种推进方式决定了没人能只守住自己最初认领的那一块。\textbf{五个人在四个版本里做过的事几乎都不重样，也几乎都超出了立项当天的分工。基于各自的实际投入，小组一致决定不进行分级评分，申请全员满分}。具体评分与理由如下。

\section{个人评分与理由}

\subsection{罗晨菲}

负责全局模块与项目统筹。V1 阶段她是工作量最集中的一个人：仓库的文件接口由她建立，页面骨架、主菜单、导航与存档体系也在这一代从零搭起，9 月 5 日 V1 检验通过时，游戏已经可以连续游玩。V2 阶段转向全局状态与正式集成，交付了项目的全局接口交接文档，并主持了界面切换规则的合并；这一代的 PRD 也由她撰写，模块边界怎么划、subPR 怎么拆，是在这一版定下来的。V3 阶段完成全局模块的对接移交。V4 阶段回到全局状态协调，负责素材、视频与音效的接入收口。开发之外，宣传视频的剪辑和过场视频的接入也由她承担。

\medskip
\noindent\textbf{得分：100}

\subsection{杨梦}

担任团队 CIO，负责美术素材与视觉规范，同时承担进度管控、会议纪要和文档整理。全篇的场景与人物素材出自她手，包括村口、祠堂、老宅三套场景的结构草图与标记草图，以及祠堂亮灯、灭灯、白灯亮起三种状态各自独立的场景图。全站视觉规范里的色板、雨、灯、字体四项由她参与制定，18 个场景能始终保持同一套明暗与色彩，靠的就是这四条。

她的工作范围在四个版本里一再外扩。V1 从零学 HTML 和 CSS，独立完成游戏入口封面页；V2 确定全站的字体与配色体系；V3 在团队人手最紧的时候补上剧情内容，逐份阅读文档之后接手旁白文本与结局文案，从美术跨到了叙事；V4 一个人做完五个结局的转场视频，又自学音频接入，把 5 段背景音乐与 7 处音效整合进游戏流程。

\medskip
\noindent\textbf{得分：100}

\subsection{于知让}

《借灯》的故事由他创作，9 月 2 日小组投票通过的正是这个剧本。以 PR 驱动需求迭代、subPR 闭环的开发方式也由他在立项当天提出，四个版本都沿这条链路推进，中途没有改过；V1 与 V3 两代的需求文档和模块分工安排由他编写和确定。

模块方面，剧情一直在他手上。V1 从零实现剧情引擎，包括节点推进、动作提交与状态记录，并定下内容与引擎分离的结构。这个结构在 V4 得到验证：46 处文本修订全部落在数据文件上，引擎代码一行没动。V2 负责剧情数据与展示适配，把阅读完成与操作提交接到统一接口。V3 定稿 28 个剧情节点的清单和五个结局的分支条件。V4 完成 46 处剧情文本修订，覆盖序幕、中段到结局，同时涉及大量的NPC对话修缮，剧情连贯性推定。

\medskip
\noindent\textbf{得分：100}

\subsection{高冰轩}

全组提交量最多的成员，累计 31 个合并请求。承担小游戏与游戏页两大块，V3 又兼了全局协调与状态存档。V1 制作项目介绍、团队与成员页面，并完成第一个小游戏；V2 集中处理正式游戏页，统一阅读组件、详情卡片、反馈提示与窄屏排版；V3 从局部界面走到跨模块联调，接入正式剧情入口与全局状态，逐条核对探索热点、小游戏命令和剧情节点能不能对上；V4 把三个新小游戏从「只有入口和结果交接」的占位状态，做成玩家真正能操作、能结束、能回到主线的完整玩法。

\medskip
\noindent\textbf{得分：100}

\subsection{卢正松}

负责探索与交互，是贯穿四个版本的联调负责人。V1 承担场景探索、NPC 对话、背包与成就。V2 定下“光源点击”这个核心交互，跑通热点与坐标的技术链路，并把同一套坐标系推给所有人对齐——这一步不做，五个人的模块就挂不到同一张场景图上。V3 把探索从局部铺开到全图，同时负责审查队友的 PR 再合入主线，保证合入节奏不断档。V4 重新设计对话框，把原本混在一起的系统对话拆成旁白、独白、人物对话三种类型，并完成物品与线索详情卡片、背包与成就页的更新。

\medskip
\noindent\textbf{得分：100}

\section{关于统一评分的说明}

\subsection{四个版本首尾相接，没有松懈的间歇}

四条版本区间是重叠的：V1 是 9.3--9.8，V2 是 9.7--9.11，V3 是 9.9--9.15，V4 是 9.16--9.20。上一代还在收尾，下一代的需求评审已经开完，十八天里没有哪一段是空的。

压力集中在后两代。V3 只有七天，要新增 28 个运行时剧情节点、四个剧情阶段、三个小游戏，还要把五个结局全部跑通；V4 只有五天，要落地 46 处剧情改写、11 处转场效果、7 段过场视频，同时把三个新小游戏补成完整玩法。V3 那段时间，组里几个人是昼夜颠倒过来的，改到凌晨两三点是常态，第二天照常开站会同步进度，工作量完全值得满分。

\subsection{每个人都身兼数职，且不止于开发}

五个人最初认领的都是一个模块，实际做下来的远不止这些。罗晨菲在全局开发之外写需求文档、剪宣传视频；杨梦从美术做到剧情、再做到音频接入；卢正松从探索模块做到全局联调，并且长期承担队友 PR 的审查；高冰轩从介绍页做到小游戏引擎，V3 又接下全局协调与状态存档；于知让从剧情引擎做到两代需求文档和开发方式的制定。全组的会议纪要、每周 DDL 清单、个人日志汇总、课堂汇报材料，也都由成员在开发之外分摊完成。

\subsection{模块互相咬合，任何一环停下都会卡住整条链路}

剧情引擎产出的节点与调查状态要交给探索模块；小游戏判定完的结果要交回统一状态与存档接口，再由剧情模块决定后续走向；场景热点挂在剧情模块给出的对象表格和 ID 清单上，ID 对不上，热点就点不中；转场和视频不绑节点编号，绑在台词文本上，读到约定的句子才播放，改写剧情的人和做转场的人必须对着同一份文本。这四条依赖链上的任何一处停摆，下游的人当天就无法往下做。也正因为这样，没有人能脱离其他人的进度单独完成任务。

\subsection{工作量有据可查}

仓库里的 229 次提交和 73 个合并请求可以逐条查证。四个版本每一代的 PRD、subPR 清单和合并记录都完整保留，谁在哪一代做了技术审查、谁做了需求验收，都能查到。每一代结束时都产出了一个可运行、可演示、可验收的增量，V4 交付的是从入口页走到五个结局、全流程可玩的完整作品。

\medskip

综上，《借灯》是五个人在十八天里连续滚动推出来的结果，每一代都站在上一代的交付之上，每个人的工作在整条链路里都无法替换。我们申请不进行分级评分，\textbf{全员按满分计。望老师批准。}
